\NativeSlot{40}{120}{%
% p068_s01

\ExamQuestion{3（2）}{5}%
先将下面三角形绕B点逆时针旋转\(90^\circ\)后画出图形，再将长方形绕D点顺时针旋转\(90^\circ\)后画出图形，最后请画出组合图形的对称轴。\par
\BookFigure{figures/p068_s01-01.tex}
}
\NativeSlot{241}{73}{%
\par\noindent%
\ExamTile{.54}{%
% p068_s02

\ExamQuestion{（3）}{5}%
移一移，画一画。图形①先向右平移\mbox{（\hspace{1.8em}）}格，再向\mbox{（\hspace{1.8em}）}平移\mbox{（\hspace{1.8em}）}格，就会与图形②拼成一个长方形。\par

% p068_s03

\noindent 将图形①移入图形②的凹口拼成长方形，把拼成的图形绕O点逆时针旋转\(90^\circ\)，请在网格图中画出旋转后的图形。\par
}\hfill%
\ExamTile{.435}{%
\BookFigure{figures/p068_s02-01.tex}
}\par
}
\NativeSlot{369}{148}{%
% p068_s04

\ExamQuestion{4.}{5}%
认真观察，填一填。\par\noindent （1）图①是以点\mbox{（\hspace{1.8em}）}为中心\mbox{（\hspace{1.8em}）}时针旋转的，在图中相应的括号里标出点A的对应点\(A^{\prime}\)。\par

% p068_s05

\ExamQuestion{（2）}{5}%
图②是以点\mbox{（\hspace{1.8em}）}为中心\mbox{（\hspace{1.8em}）}时针旋转的，在图中相应的括号里标出点B的对应点\(B^{\prime}\)。\par
\BookFigure{figures/p068_s04-01.tex}
}
\NativeSlot{619}{60}{%
\ExamGroup{四、动脑筋想一想。{\normalfont\kaishu\fontsize{10.8bp}{14bp}\selectfont （共2分）}}

% p068_s06

\ExamQuestion{操作}{7}%
将一个三角形先向右平移6格，再绕中心点顺时针旋转\(90^\circ\)，得到现在的图形。如果想让它回到原来的位置，应该怎样操作？请写完整步骤。\par
}
